% Complete Section 4.5 and Box 4.1, original PDF 54-60 / printed 44-50.
% The printed signs and the last summand in Box 4.1 are retained, not corrected.
% block: hodge-r3-heading page:54 type:section
\subsection{\BT{Example: The Hodge Laplacian in $\mathbb R^3$}{Example: $\mathbb R^3$ 中的 Hodge Laplacian}}\label{sec:hodge-r3}

% block: hodge-r3-introduction pages:54,55 type:ordinary-paragraph
\BT{We now return to our running example, the de Rham complex on a domain with Lipschitz boundary in $\mathbb R^3$. We saw in Section~\ref{sec:l2-de-rham-complex} that this forms a closed Hilbert complex. Therefore, for each $k=0,1,2,3$, we obtain a Hodge Laplace problem associated to the complex, which is a well-posed boundary value problem in partial differential equations. To interpret these problems we just have to interpret the spaces $V^k$ and $V_k^*$ and the operators $d^k$ and $d_k^*$, for which we will have frequent recourse to the domain complex \eqref{eq:l2-de-rham-domain-complex} and the dual domain complex \eqref{eq:l2-dual-domain-complex}.}{现在回到我们贯穿使用的例子, 即 $\mathbb R^3$ 中具有 Lipschitz 边界的区域上的 de Rham 复形. 在 Section~\ref{sec:l2-de-rham-complex} 中已看到, 它构成闭 Hilbert 复形. 因此, 对每个 $k=0,1,2,3$, 都得到与该复形相伴的 Hodge Laplace 问题, 它是一个适定的偏微分方程边值问题. 为了解释这些问题, 只需解释空间 $V^k$ 和 $V_k^*$ 以及算子 $d^k$ 和 $d_k^*$; 为此我们将频繁使用定义域复形 \eqref{eq:l2-de-rham-domain-complex} 及对偶定义域复形 \eqref{eq:l2-dual-domain-complex}.}

% block: hodge-r3-k0-heading page:55 type:subsection
\subsubsection{$k=0$}\label{subsec:hodge-r3-k0}

% block: hodge-r3-k0-domain page:55 type:ordinary-paragraph
\BT{Assuming that the domain is connected, the space $\mathsf H^0$ of harmonic forms consists only of the constants and $P_{\mathsf H}f=\bar f$, the mean value of $f$ on $\Omega$. Since $d_0^*=0$, $L^0u=(dd^*+d^*d)u=d^*du=-\operatorname{div}\operatorname{grad}u=-\Delta u$, with domain}{假设区域连通, 则调和形式空间 $\mathsf H^0$ 仅由常数构成, 且 $P_{\mathsf H}f=\bar f$, 即 $f$ 在 $\Omega$ 上的平均值. 由于 $d_0^*=0$, 有 $L^0u=(dd^*+d^*d)u=d^*du=-\operatorname{div}\operatorname{grad}u=-\Delta u$, 定义域为}
\[
\begin{aligned}
D(L)&=\{\,u\in H^1(\Omega)\mid\operatorname{grad}u\in\mathring H(\operatorname{div},\Omega)\,\}\\
&=\{\,u\in H^1(\Omega)\mid\Delta u\in L^2(\Omega),\ \partial u/\partial n=0\text{ on }\partial\Omega\,\},
\end{aligned}
\]
\BT{where $\Delta$ denotes the ordinary Laplacian $\Delta=\sum_{i=1}^3\partial^2/\partial x_i^2$. Given $f\in L^2(\Omega)$, the strong formulation \eqref{eq:hodge-strong-formulation} then seeks $u\in H^1$ such that}{其中 $\Delta$ 表示通常的 Laplacian $\Delta=\sum_{i=1}^3\partial^2/\partial x_i^2$. 给定 $f\in L^2(\Omega)$, 强表述 \eqref{eq:hodge-strong-formulation} 就是求 $u\in H^1$, 使得}
\begin{equation}\label{eq:hodge-r3-neumann-strong}
-\Delta u=f-\bar f\text{ in }\Omega,\qquad\frac{\partial u}{\partial n}=0\text{ on }\partial\Omega,\qquad\int u=0.
\end{equation}

% block: hodge-r3-neumann-interpretation page:55 type:ordinary-paragraph
\BT{Thus the $k=0$ Hodge Laplace problem associated to the de Rham complex is just the usual Neumann problem for Poisson's equation, the adjustment by $\bar f$ being required to make the data $f$ compatible with the homogeneous Neumann boundary condition.}{因此, 与 de Rham 复形相伴的 $k=0$ Hodge Laplace 问题, 就是 Poisson 方程通常的 Neumann 问题; 减去 $\bar f$ 的调整是为了使数据 $f$ 与齐次 Neumann 边界条件相容.}

% block: hodge-r3-neumann-weak page:55 type:ordinary-paragraph
\BT{Next we turn to the primal and mixed weak formulations. Since $V^{-1}=0$, for $k=0$, $\sigma$ and $\tau$ vanish in the mixed weak formulation \eqref{eq:hodge-mixed-formulation}. Thus there is no distinction between the two weak formulations for $k=0$. They both seek $u\in H^1$ such that}{接下来讨论原始和混合弱表述. 由于 $V^{-1}=0$, 当 $k=0$ 时, 混合弱表述 \eqref{eq:hodge-mixed-formulation} 中的 $\sigma$ 和 $\tau$ 均消失. 因此在 $k=0$ 时两种弱表述没有区别. 它们都是求 $u\in H^1$, 使得}
\begin{equation}\label{eq:hodge-r3-neumann-weak}
\int\operatorname{grad}u\cdot\operatorname{grad}v=\int(f-\bar f)v,\qquad v\in H^1,\qquad\int_\Omega u=0.
\end{equation}

% block: hodge-r3-neumann-equivalence page:55 type:ordinary-paragraph
\BT{We have seen in Theorem~\ref{thm:hodge-formulation-equivalence} that this problem is equivalent to the strongly formulated problem \eqref{eq:hodge-r3-neumann-strong}. For concreteness, we now check directly that \eqref{eq:hodge-r3-neumann-weak} implies \eqref{eq:hodge-r3-neumann-strong}, the other direction being easier. Taking $v\in C_0^\infty(\Omega)$ in \eqref{eq:hodge-r3-neumann-weak} we see that}{由 Theorem~\ref{thm:hodge-formulation-equivalence} 已知, 此问题等价于强表述问题 \eqref{eq:hodge-r3-neumann-strong}. 为使讨论具体, 现在直接验证 \eqref{eq:hodge-r3-neumann-weak} 蕴含 \eqref{eq:hodge-r3-neumann-strong}, 另一个方向更容易. 在 \eqref{eq:hodge-r3-neumann-weak} 中取 $v\in C_0^\infty(\Omega)$, 可见}
\begin{equation}\label{eq:hodge-r3-neumann-distribution}
-\operatorname{div}\operatorname{grad}u=f-\bar f
\end{equation}
\BT{in the sense of distributions and so, since $f\in L^2$, $\operatorname{grad}u$ belongs to $H(\operatorname{div})$. We may therefore apply the integration-by-parts formula \eqref{eq:normal-trace-green} with $w=\operatorname{grad}u$ and then apply \eqref{eq:hodge-r3-neumann-weak} and \eqref{eq:hodge-r3-neumann-distribution} to get}{在分布意义下成立, 因此由 $f\in L^2$ 可知, $\operatorname{grad}u$ 属于 $H(\operatorname{div})$. 故可在分部积分公式 \eqref{eq:normal-trace-green} 中取 $w=\operatorname{grad}u$, 再利用 \eqref{eq:hodge-r3-neumann-weak} 和 \eqref{eq:hodge-r3-neumann-distribution}, 得到}
\[
\langle\gamma_n(\operatorname{grad}u),\operatorname{tr}v\rangle=\int(\operatorname{div}\operatorname{grad}u)v+\int\operatorname{grad}u\cdot\operatorname{grad}v=0,\quad v\in H^1.
\]
\BT{Since $\operatorname{tr}v$ can be an arbitrary element of $H^{1/2}(\partial\Omega)$ we conclude that $\gamma_n(\operatorname{grad}u)=0$, i.e., $\partial u/\partial n=0$ on the boundary. Notice that, in the weak formulation, the boundary condition is imposed \emph{naturally}, meaning that it is implied by the equations via integration by parts rather than being imposed in the space in which $u$ is sought.}{由于 $\operatorname{tr}v$ 可以是 $H^{1/2}(\partial\Omega)$ 的任意元素, 可知 $\gamma_n(\operatorname{grad}u)=0$, 即边界上有 $\partial u/\partial n=0$. 注意, 在弱表述中, 边界条件是\emph{自然}施加的, 即通过分部积分由方程蕴含, 而不是施加在求解 $u$ 所在的空间中.}

% block: hodge-r3-k0-boundary-problems page:55 type:ordinary-paragraph
\BT{Turning to the $\mathcal B$ and $\mathcal B^*$ problems, we do not find anything new, since $\mathcal B^0=0$. Therefore, the $\mathcal B$ problem does not arise for $k=0$, while the $\mathcal B^*$ problem is essentially the full Hodge Laplacian.}{再看 $\mathcal B$ 和 $\mathcal B^*$ 问题, 由于 $\mathcal B^0=0$, 并没有新的内容. 因此, 在 $k=0$ 时不存在 $\mathcal B$ 问题, 而 $\mathcal B^*$ 问题本质上就是完整的 Hodge Laplacian 问题.}

% block: hodge-r3-k0-summary page:56 type:ordinary-paragraph
\BT{Summarizing the situation for $k=0$, the Hodge Laplacian for $0$-forms is just the Neumann problem for the Poisson equation. The primal and mixed weak formulations coincide in this case, being just the standard weak formulation for the Neumann problem, in which the boundary condition is imposed \emph{naturally}.}{概括 $k=0$ 的情形, $0$-形式的 Hodge Laplacian 就是 Poisson 方程的 Neumann 问题. 此时原始和混合弱表述一致, 都是 Neumann 问题的标准弱表述, 其中边界条件是\emph{自然}施加的.}

% block: hodge-r3-k1-heading page:56 type:subsection
\subsubsection{$k=1$}\label{subsec:hodge-r3-k1}

% block: hodge-r3-k1-strong page:56 type:ordinary-paragraph
\BT{The case $k=1$ is more interesting. The space of harmonic forms $\mathsf H^1$ is a space of vector fields whose dimension is $b_1$, the first Betti number. The strong formulation seeks $u\in H(\operatorname{curl})\cap\mathring H(\operatorname{div})$ with $\operatorname{curl}u\in\mathring H(\operatorname{curl})$ and $\operatorname{div}u\in H^1$ satisfying the partial differential equation}{$k=1$ 的情形更有意思. 调和形式空间 $\mathsf H^1$ 是向量场空间, 其维数为 $b_1$, 即第一个 Betti 数. 强表述是求 $u\in H(\operatorname{curl})\cap\mathring H(\operatorname{div})$, 满足 $\operatorname{curl}u\in\mathring H(\operatorname{curl})$, $\operatorname{div}u\in H^1$, 以及偏微分方程}
\begin{equation}\label{eq:hodge-r3-vector-laplace}
\operatorname{curl}\operatorname{curl}u-\operatorname{grad}\operatorname{div}u=f-P_{\mathsf H}f\text{ in }\Omega.
\end{equation}
\BT{Note that the differential operator $\operatorname{curl}\operatorname{curl}-\operatorname{grad}\operatorname{div}$ is the negative vector Laplacian $-\vec\Delta$ in that \[\operatorname{curl}\operatorname{curl}u-\operatorname{grad}\operatorname{div}u=(-\Delta u_1,-\Delta u_2,-\Delta u_3)\] for a smooth vector field $u=(u_1,u_2,u_3)$.}{注意, 微分算子 $\operatorname{curl}\operatorname{curl}-\operatorname{grad}\operatorname{div}$ 是负的向量 Laplacian $-\vec\Delta$, 因为对光滑向量场有 \[\operatorname{curl}\operatorname{curl}u-\operatorname{grad}\operatorname{div}u=(-\Delta u_1,-\Delta u_2,-\Delta u_3)\] 其中 $u=(u_1,u_2,u_3)$.}

% block: hodge-r3-magnetic-boundary page:56 type:ordinary-paragraph
\BT{The definition of the domain $D(L^1)$ implies the boundary conditions}{定义域 $D(L^1)$ 的定义蕴含边界条件}
\begin{equation}\label{eq:hodge-r3-magnetic-boundary}
u\cdot n=0,\qquad(\operatorname{curl}u)\times n=0\text{ on }\partial\Omega,
\end{equation}
\BT{and, of course, we have the side condition $u\perp\mathsf H^1$ to impose uniqueness. The conditions \eqref{eq:hodge-r3-magnetic-boundary} are sometimes referred to as \emph{magnetic boundary conditions}, for reasons that will become clearer when we discuss Maxwell's equations in Section~\ref{sec:maxwell-equations}.}{当然, 还需要附加条件 $u\perp\mathsf H^1$ 来保证唯一性. 条件 \eqref{eq:hodge-r3-magnetic-boundary} 有时称为\emph{磁边界条件}, 其原因将在 Section~\ref{sec:maxwell-equations} 讨论 Maxwell 方程时变得更清楚.}

% block: hodge-r3-k1-primal page:56 type:ordinary-paragraph
\BT{The primal weak formulation seeks $u\in H(\operatorname{curl})\cap\mathring H(\operatorname{div})$ with}{原始弱表述是求 $u\in H(\operatorname{curl})\cap\mathring H(\operatorname{div})$, 使得}
\begin{equation}\label{eq:hodge-r3-magnetic-primal}
\begin{gathered}
\langle\operatorname{curl}u,\operatorname{curl}v\rangle+\langle\operatorname{div}u,\operatorname{div}v\rangle=\langle f-P_{\mathsf H}f,v\rangle,\\
v\in H(\operatorname{curl})\cap\mathring H(\operatorname{div}).
\end{gathered}
\end{equation}
\BT{Notice that, in this formulation, the boundary condition $u\cdot n=0$ is essential (imposed in the space where $u$ is sought), while the condition $\operatorname{curl}u\times n=0$ is natural.}{注意, 在此表述中, 边界条件 $u\cdot n=0$ 是本质的 (施加在求解 $u$ 所在的空间中), 而条件 $\operatorname{curl}u\times n=0$ 是自然的.}

% block: hodge-r3-k1-mixed page:56 type:ordinary-paragraph
\BT{By contrast, the mixed weak formulation seeks $\sigma\in H^1$, $u\in H(\operatorname{curl})$, $p\in\mathsf H^1$, without any imposed boundary conditions, satisfying}{相比之下, 混合弱表述是求 $\sigma\in H^1$, $u\in H(\operatorname{curl})$, $p\in\mathsf H^1$, 不施加任何边界条件, 并满足}
\begin{equation}\label{eq:hodge-r3-magnetic-mixed}
\begin{aligned}
\int\sigma\tau-\int u\cdot\operatorname{grad}\tau&=0,&\tau&\in H^1,\\
\int\operatorname{grad}\sigma\cdot v+\int\operatorname{curl}u\cdot\operatorname{curl}v+\int p\cdot v&=\int f\cdot v,&v&\in H(\operatorname{curl}),\\
\int u\cdot q&=0,&q&\in\mathsf H^1.
\end{aligned}
\end{equation}
\BT{The boundary conditions \eqref{eq:hodge-r3-magnetic-boundary} are both natural in this formulation, arising from the integration by parts in the first and second equations, respectively. Notice that, in addition to the unknown vector field $u$, this formulation takes as a second unknown the scalar function $\sigma=-\operatorname{div}u$ (and of course $p=P_{\mathsf H}f$).}{在此表述中, 边界条件 \eqref{eq:hodge-r3-magnetic-boundary} 都是自然的, 分别来自第一个和第二个方程中的分部积分. 注意, 除未知向量场 $u$ 外, 此表述还以标量函数 $\sigma=-\operatorname{div}u$ 为第二个未知量 (当然还有 $p=P_{\mathsf H}f$).}

% block: hodge-r3-k1-b-problem page:56 type:ordinary-paragraph
\BT{In the case $f\in\mathcal B^1$, i.e., $f$ is a gradient, then the solution $u$ is a gradient as well, which we can express by saying that $\operatorname{curl}u=0$ and $u\perp\mathsf H^1$. Thus the $\mathcal B^1$ problem is to find $u$ such that}{当 $f\in\mathcal B^1$, 即 $f$ 为梯度时, 解 $u$ 也是梯度, 这可以表述为 $\operatorname{curl}u=0$ 且 $u\perp\mathsf H^1$. 因此 $\mathcal B^1$ 问题是求 $u$, 使得}
\[
-\operatorname{grad}\operatorname{div}u=f,\ \operatorname{curl}u=0\text{ in }\Omega,\qquad u\cdot n=0\text{ on }\partial\Omega,\qquad u\perp\mathsf H^1.
\]

% block: hodge-r3-div-curl-application page:57 type:ordinary-paragraph
\BT{From the abstract theory, we know that this problem has a unique solution for $f\in\mathcal B^1$ which, moreover, can be computed by solving the full Hodge Laplacian problem for that same $f$ (e.g., via the mixed weak formulation). One application of such a problem is toward solving the div-curl problem}{由抽象理论可知, 对 $f\in\mathcal B^1$, 此问题具有唯一解; 而且, 可通过对同一 $f$ 求解完整的 Hodge Laplacian 问题 (例如通过混合弱表述) 来计算此解. 此类问题的一个应用是求解 div-curl 问题}
\[
\operatorname{div}u=F,\qquad\operatorname{curl}u=0,
\]
\BT{which can be done by taking $f=\operatorname{grad}F$ (with $\int F=0$ required for compatibility with the boundary condition on $u$).}{这可以通过取 $f=\operatorname{grad}F$ 实现 (为与 $u$ 的边界条件相容, 需要 $\int F=0$).}

% block: hodge-r3-k1-bstar-problem page:57 type:ordinary-paragraph
\BT{On the other hand, the $\mathcal B^*$ problem arises when $f\in\mathcal B_1^*$, meaning $f\in\operatorname{curl}\mathring H(\operatorname{curl})$. Then $u$ is determined by}{另一方面, 当 $f\in\mathcal B_1^*$, 即 $f\in\operatorname{curl}\mathring H(\operatorname{curl})$ 时, 就得到 $\mathcal B^*$ 问题. 此时 $u$ 由下式确定:}
\[
\begin{gathered}
\operatorname{curl}\operatorname{curl}u=f,\ \operatorname{div}u=0\text{ in }\Omega,\\
u\cdot n=0,\ \operatorname{curl}u\times n=0\text{ on }\partial\Omega,\qquad u\perp\mathsf H^1.
\end{gathered}
\]
\BT{Such curl-curl problems arise in electromagnetics. Again, the solution may be determined by solving the full Hodge Laplace problem.}{这种 curl-curl 问题出现在电磁学中. 同样, 可通过求解完整的 Hodge Laplace 问题来确定其解.}

% block: hodge-r3-k2-heading page:57 type:subsection
\subsubsection{$k=2$}\label{subsec:hodge-r3-k2}

% block: hodge-r3-k2-domain page:57 type:ordinary-paragraph
\BT{This case again leads to the vector Laplace equation \eqref{eq:hodge-r3-vector-laplace}, but with a different set of boundary conditions (and different harmonic forms with $\dim\mathsf H^2=b_2$). The domain of $L^2$ is the set of $u\in H(\operatorname{div})\cap\mathring H(\operatorname{curl})$ with $\operatorname{div}u\in\mathring H^1$ and $\operatorname{curl}u\in H(\operatorname{curl})$. Thus we obtain \emph{electric boundary conditions}:}{此情形仍得到向量 Laplace 方程 \eqref{eq:hodge-r3-vector-laplace}, 但边界条件不同 (调和形式也不同, 其中 $\dim\mathsf H^2=b_2$). $L^2$ 的定义域由满足 $\operatorname{div}u\in\mathring H^1$ 和 $\operatorname{curl}u\in H(\operatorname{curl})$ 的 $u\in H(\operatorname{div})\cap\mathring H(\operatorname{curl})$ 构成. 因此得到\emph{电边界条件}:}
\[
u\times n=0,\ \operatorname{div}u=0\text{ on }\partial\Omega.
\]

% block: hodge-r3-k2-mixed page:57 type:ordinary-paragraph
\BT{The first of these conditions is essential in the primal variational formulation, while the second is natural. Both are, of course, natural in the mixed weak formulation, which seeks $\sigma\in H(\operatorname{curl})$, $u\in H(\operatorname{div})$, $p\in\mathsf H^2$ such that}{这些条件中, 第一个在原始变分表述中是本质的, 第二个是自然的. 当然, 在混合弱表述中两者都是自然的; 该表述是求 $\sigma\in H(\operatorname{curl})$, $u\in H(\operatorname{div})$, $p\in\mathsf H^2$, 使得}
\begin{equation}\label{eq:hodge-r3-electric-mixed}
\begin{aligned}
\int\sigma\cdot\tau-\int u\cdot\operatorname{curl}\tau&=0,&\tau&\in H(\operatorname{curl}),\\
\int\operatorname{curl}\sigma\cdot v+\int\operatorname{div}u\operatorname{div}v+\int p\cdot v&=\int f\cdot v,&v&\in H(\operatorname{div}),\\
\int u\cdot q&=0,&q&\in\mathsf H^2.
\end{aligned}
\end{equation}
\BT{Notice that this formulation is very different from the mixed weak formulation \eqref{eq:hodge-r3-magnetic-mixed} for $k=1$. Here the additional field $\sigma$ is a vector field, while it was a scalar field for $k=1$.}{注意, 此表述与 $k=1$ 时的混合弱表述 \eqref{eq:hodge-r3-magnetic-mixed} 很不相同. 这里的附加场 $\sigma$ 是向量场, 而在 $k=1$ 时它是标量场.}

% block: hodge-r3-k2-boundary-problems page:57 type:ordinary-paragraph
\BT{The $\mathcal B$ problem for $k=2$ is}{当 $k=2$ 时, $\mathcal B$ 问题为}
\[
\operatorname{curl}\operatorname{curl}u=f,\ \operatorname{div}u=0\text{ in }\Omega,\qquad u\times n=0,\text{ on }\partial\Omega,\qquad u\perp\mathsf H^2,
\]
\BT{which has a unique solution if $f\in\operatorname{curl}H(\operatorname{curl})$. The $\mathcal B_2^*$ problem is}{若 $f\in\operatorname{curl}H(\operatorname{curl})$, 则它具有唯一解. $\mathcal B_2^*$ 问题为}
\[
\begin{gathered}
-\operatorname{grad}\operatorname{div}u=f,\ \operatorname{curl}u=0\text{ in }\Omega,\\
u\times n=0,\ \operatorname{div}u=0\text{ on }\partial\Omega,\qquad u\perp\mathsf H^2,
\end{gathered}
\]
\BT{solvable for $f\in\operatorname{grad}\mathring H^1$. Notice that the differential equations, but not the boundary conditions, for the $\mathcal B^2$ problem coincide with those of the $\mathcal B_1^*$ problem and those for the $\mathcal B_2^*$ problem coincide with those of the $\mathcal B^1$ problem.}{对 $f\in\operatorname{grad}\mathring H^1$ 可解. 注意, $\mathcal B^2$ 问题的微分方程与 $\mathcal B_1^*$ 问题的微分方程一致, $\mathcal B_2^*$ 问题的微分方程与 $\mathcal B^1$ 问题的微分方程一致, 但相应的边界条件并不相同.}

% block: hodge-r3-table page:58 type:table
\begin{table}[htbp]\centering\small
\caption{\BT{The Hodge Laplace problem for the de Rham complex. $\sigma$ is the extra variable in mixed weak formulation. Underlined boundary conditions (the ones not involving derivatives) are essential in the primal weak formulation, while the others are natural. All the boundary conditions are natural in the mixed weak formulation.}{de Rham 复形的 Hodge Laplace 问题. $\sigma$ 是混合弱表述中的附加变量. 带下划线的边界条件 (不涉及导数的那些条件) 在原始弱表述中是本质的, 其余是自然的. 所有边界条件在混合弱表述中都是自然的.}}\label{tab:hodge-r3-boundary-conditions}
\setlength{\tabcolsep}{4pt}
\begin{tabular}{cccll}\hline
$k$&$L^k$&$\sigma$&\BT{Boundary conditions}{边界条件}&\BT{Harmonic forms}{调和形式}\\\hline
$0$&$-\Delta$&$0$&Neumann: $\partial u/\partial n=0$&$\mathsf H^0$ $(\dim b_0)$\\\noalign{\vskip4pt}
$1$&$-\vec\Delta$&$-\operatorname{div}u$&\BT{magnetic}{磁}: $\underline{u\cdot n=0}$, $\operatorname{curl}u\times n=0$&$\mathsf H^1$ $(\dim b_1)$\\\noalign{\vskip4pt}
$2$&$-\vec\Delta$&$\operatorname{curl}u$&\BT{electric}{电}: $\underline{u\times n=0}$, $\operatorname{div}u=0$&$\mathsf H^2$ $(\dim b_2)$\\\noalign{\vskip4pt}
$3$&$-\Delta$&$-\operatorname{grad}u$&Dirichlet: $\underline{u=0}$&$\mathsf H^3=0$\\\hline
\end{tabular}
\end{table}

% block: hodge-r3-k3-heading page:58 type:subsection
\subsubsection{$k=3$}\label{subsec:hodge-r3-k3}

% block: hodge-r3-k3-dirichlet page:58 type:ordinary-paragraph
\BT{Both $\mathcal B_3^*$ and $\mathsf H^3$ vanish. The domain of the Laplacian $D(L^3)=\{\,u\in\mathring H^1\mid\Delta u\in L^2\,\}$, so the boundary value problem is just the Dirichlet problem for Poisson's equation:}{空间 $\mathcal B_3^*$ 和 $\mathsf H^3$ 都为零. Laplacian 的定义域为 $D(L^3)=\{\,u\in\mathring H^1\mid\Delta u\in L^2\,\}$, 因此该边值问题就是 Poisson 方程的 Dirichlet 问题:}
\[
-\Delta u=f\text{ in }\Omega,\qquad u=0\text{ on }\partial\Omega.
\]
\BT{The primal weak formulation is the obvious one: find $u\in\mathring H^1$ such that}{其原始弱表述是显然的: 求 $u\in\mathring H^1$, 使得}
\[
\int\operatorname{grad}u\cdot\operatorname{grad}v=\int fv,\qquad v\in\mathring H^1.
\]
\BT{The mixed weak formulation seeks $\sigma\in H(\operatorname{div})$, $u\in L^2$ such that}{混合弱表述是求 $\sigma\in H(\operatorname{div})$, $u\in L^2$, 使得}
\[
\begin{gathered}
\int\sigma\cdot\tau-\int u\operatorname{div}\tau=0,\qquad\tau\in H(\operatorname{div}),\\
\int\operatorname{div}\sigma\,v=\int fv,\qquad v\in L^2.
\end{gathered}
\]
\BT{The Dirichlet condition is natural in the mixed formulation.}{在混合表述中, Dirichlet 条件是自然的.}

% block: hodge-r3-table-summary page:58 type:ordinary-paragraph
\BT{Table~\ref{tab:hodge-r3-boundary-conditions} collects key features of the Hodge Laplacian problems for the de Rham complex.}{Table~\ref{tab:hodge-r3-boundary-conditions} 汇集了 de Rham 复形的 Hodge Laplacian 问题的主要特征.}

% block: hodge-r3-boundary-duality-heading page:58 type:subsection
\subsubsection{\BT{The de Rham complex with boundary conditions and duality}{带边界条件的 de Rham 复形与对偶性}}\label{subsec:de-rham-boundary-duality}

% block: hodge-r3-boundary-complex page:58 type:ordinary-paragraph
\BT{We have just considered the application of the abstract theory to the $L^2$ de Rham complex, whose domain complex is given by \eqref{eq:l2-de-rham-domain-complex}. Instead, we may apply the theory to the \emph{$L^2$ de Rham complex with boundary conditions}, by which we mean the Hilbert complex whose base spaces remain $L^2$ but whose domain complex is}{刚才讨论了抽象理论在 $L^2$ de Rham 复形中的应用, 其定义域复形由 \eqref{eq:l2-de-rham-domain-complex} 给出. 也可以将理论应用于\emph{带边界条件的 $L^2$ de Rham 复形}, 即基空间仍为 $L^2$, 而定义域复形为下式的 Hilbert 复形:}
\[
0\to\mathring H^1\xrightarrow{\operatorname{grad}}\mathring H(\operatorname{curl})\xrightarrow{\operatorname{curl}}\mathring H(\operatorname{div})\xrightarrow{\operatorname{div}}L^2\to0.
\]
\BT{This Hilbert complex is essentially the dual complex to \eqref{eq:l2-de-rham-domain-complex}, i.e., \eqref{eq:l2-dual-domain-complex}, except for some signs and for a change of indexing. Specifically, the index $k$ space of the de Rham complex with boundary conditions is the index $3-k$ space of the dual complex to the original de Rham complex. The dual complex to the de Rham complex with boundary conditions is then}{除若干符号和指标的变化外, 此 Hilbert 复形本质上就是 \eqref{eq:l2-de-rham-domain-complex} 的对偶复形, 即 \eqref{eq:l2-dual-domain-complex}. 具体地, 带边界条件的 de Rham 复形中指标为 $k$ 的空间, 就是原 de Rham 复形的对偶复形中指标为 $3-k$ 的空间. 因此, 带边界条件的 de Rham 复形的对偶复形为}
\[
0\leftarrow L^2\xleftarrow{-\operatorname{div}}H(\operatorname{div})\xleftarrow{\operatorname{curl}}H(\operatorname{curl})\xleftarrow{-\operatorname{grad}}H^1\leftarrow0,
\]
\BT{i.e., the original de Rham complex, except for signs and indexing.}{即除符号和指标外, 就是原 de Rham 复形.}

% block: hodge-r3-boundary-table page:59 type:table
\begin{table}[htbp]\centering\small
\caption{\BT{The Hodge Laplace problem for the de Rham complex with boundary conditions. All boundary conditions are essential for the mixed formulation and underlined boundary conditions are essential for the primal formulation. See Table~\ref{tab:hodge-r3-boundary-conditions}.}{带边界条件的 de Rham 复形的 Hodge Laplace 问题. 在混合表述中, 所有边界条件都是本质的; 在原始表述中, 带下划线的边界条件是本质的. 见 Table~\ref{tab:hodge-r3-boundary-conditions}.}}\label{tab:hodge-r3-relative-boundary-conditions}
\setlength{\tabcolsep}{3pt}
\begin{tabular}{cccll}\hline
$k$&$L^k$&$\sigma$&\BT{Boundary conditions}{边界条件}&\BT{Harmonic forms}{调和形式}\\\hline
$0$&$-\Delta$&$0$&Dirichlet: $\underline{u=0}$&$\mathring{\mathsf H}^0=0$\\\noalign{\vskip4pt}
$1$&$-\vec\Delta$&$-\operatorname{div}u$&\BT{electric}{电}: $\underline{u\times n=0}$, $\operatorname{div}u=0$&$\mathring{\mathsf H}^1=\mathsf H^2$ $(\dim b_2)$\\\noalign{\vskip4pt}
$2$&$-\vec\Delta$&$\operatorname{curl}u$&\BT{magnetic}{磁}: $\underline{u\cdot n=0}$, $\operatorname{curl}u\times n=0$&$\mathring{\mathsf H}^2=\mathsf H^1$ $(\dim b_1)$\\\noalign{\vskip4pt}
$3$&$-\Delta$&$-\operatorname{grad}u$&Neumann: $\partial u/\partial n=0$&$\mathring{\mathsf H}^3=\mathsf H^0$ $(\dim b_0)$\\\hline
\end{tabular}
\end{table}

% block: hodge-r3-poincare-lefschetz page:59 type:ordinary-paragraph
\BT{If we write $\mathring{\mathcal B}^k$, $\mathring{\mathcal Z}^k$, and $\mathring{\mathsf H}^k$ for the spaces of boundaries, cycles, and harmonic forms of the complex with boundary conditions, then $\mathring{\mathcal B}^k=\mathcal B^{3-k}$ and similarly for $\mathcal Z$ and $\mathsf H$. (The relation $\mathsf H^{n-k}=\mathring{\mathsf H}^k$ on an $n$-dimensional manifold is known as \emph{Poincaré–Lefschetz duality}. For manifolds without boundary it becomes $\mathsf H^{n-k}=\mathsf H^k$, which is Poincaré duality.) It is easy to check that the strong formulation and the primal weak formulation for the complex with boundary conditions coincide with those for the ordinary de Rham complex, except for an index change, so we obtain the same boundary value problems as above. The $k=0$ problem, for instance, is the Poisson equation with Dirichlet boundary conditions. The mixed weak formulation for the problems with $k>0$, however, brings in something new. For $k=1$, for example, we seek $\sigma\in\mathring H^1$, $u\in\mathring H(\operatorname{curl})$, $p\in\mathring{\mathsf H}^1=\mathsf H^2$ such that}{若用 $\mathring{\mathcal B}^k$, $\mathring{\mathcal Z}^k$ 和 $\mathring{\mathsf H}^k$ 表示带边界条件的复形的边界空间, 闭链空间和调和形式空间, 则有 $\mathring{\mathcal B}^k=\mathcal B^{3-k}$, 对 $\mathcal Z$ 和 $\mathsf H$ 也有类似关系. (在 $n$ 维流形上的关系 $\mathsf H^{n-k}=\mathring{\mathsf H}^k$ 称为 \emph{Poincaré–Lefschetz 对偶性}. 对无边界流形, 它成为 $\mathsf H^{n-k}=\mathsf H^k$, 即 Poincaré 对偶性.) 容易验证, 除指标变化外, 带边界条件的复形的强表述和原始弱表述与通常的 de Rham 复形的相应表述一致, 因而得到与上面相同的边值问题. 例如 $k=0$ 的问题就是带 Dirichlet 边界条件的 Poisson 方程. 然而, 对 $k>0$ 的问题, 混合弱表述带来了新的内容. 例如当 $k=1$ 时, 求 $\sigma\in\mathring H^1$, $u\in\mathring H(\operatorname{curl})$, $p\in\mathring{\mathsf H}^1=\mathsf H^2$, 使得}
\[
\begin{aligned}
\int\sigma\tau-\int u\cdot\operatorname{grad}\tau&=0,&\tau&\in\mathring H^1,\\
\int\operatorname{grad}\sigma\cdot v+\int\operatorname{curl}u\cdot\operatorname{curl}v+\int p\cdot v&=\int f\cdot v,&v&\in\mathring H(\operatorname{curl}).\\
\int u\cdot q&=0,&q&\in\mathsf H^2.
\end{aligned}
\]

% block: hodge-r3-relative-mixed-comparison page:59 type:ordinary-paragraph
\BT{This is a mixed formulation for the vector Laplacian with boundary conditions $\operatorname{div}u=0$, $u\times n=0$. These are the same boundary conditions as for the $k=2$ problem for the ordinary de Rham complex. However, the mixed method \eqref{eq:hodge-r3-electric-mixed} for that problem was altogether different, using a vector field for the auxiliary variable $\sigma$ rather than a scalar field, as here, and imposing the boundary conditions naturally, rather than essentially as here.}{这是向量 Laplacian 在边界条件 $\operatorname{div}u=0$, $u\times n=0$ 下的混合表述. 这些边界条件与通常的 de Rham 复形中 $k=2$ 问题的边界条件相同. 但是, 那个问题的混合方法 \eqref{eq:hodge-r3-electric-mixed} 完全不同, 辅助变量 $\sigma$ 使用向量场而不是这里的标量场, 并且边界条件是自然施加的, 而不是像这里这样作为本质条件施加.}

% block: hodge-r3-relative-summary page:59 type:ordinary-paragraph
\BT{A similar situation holds for $k=2$ and $k=3$, which we leave to the reader to work out in detail. A summary is presented in Table~\ref{tab:hodge-r3-relative-boundary-conditions}.}{对 $k=2$ 和 $k=3$ 也有类似情形, 细节留待读者推导. Table~\ref{tab:hodge-r3-relative-boundary-conditions} 给出了汇总.}

% block: hilbert-summary-heading page:60 type:ordinary-paragraph
\begin{tcolorbox}[enhanced,breakable,colback=KuangNei,frame hidden,boxrule=0pt,sharpish corners,before skip=10pt,after skip=10pt]
\refstepcounter{booksummary}\label{box:hilbert-complexes-summary}
\noindent\textbf{Box~\thebooksummary:} \BT{Summary of Chapter~\ref{chap:hilbert-complexes} on Hilbert complexes.}{Chapter~\ref{chap:hilbert-complexes} Hilbert 复形内容摘要.}\par

% block: hilbert-summary-definition page:60 type:ordinary-paragraph
\BT{A \emph{Hilbert complex} is a sequence of Hilbert spaces $W^k$ and closed densely defined operators $d^k$ from one to the next such that $d^k$ maps its domain into the null space of $d^{k+1}$. The Hilbert complex is called \emph{closed} if the $d^k$ all have closed range.}{\emph{Hilbert 复形}是由 Hilbert 空间 $W^k$ 以及从一个空间到下一个空间的稠定闭算子 $d^k$ 构成的序列, 其中 $d^k$ 将其定义域映入 $d^{k+1}$ 的零空间. 若所有 $d^k$ 都具有闭像, 则称此 Hilbert 复形是\emph{闭的}.}

% block: hilbert-summary-domain-complex page:60 type:ordinary-paragraph
\BT{Setting $V^k=D(d^k)$, we get the domain complex}{令 $V^k=D(d^k)$, 得到定义域复形}
\[
0\to V^0\xrightarrow{d}V^1\xrightarrow{d}\cdots\xrightarrow{d}V^n\to0,
\]
\BT{a cochain complex of Hilbert spaces and bounded linear operators.}{它是由 Hilbert 空间和有界线性算子构成的上链复形.}

% block: hilbert-summary-cocycles page:60 type:ordinary-paragraph
\BT{We denote the cocycles and coboundaries as $\mathcal Z^k$ and $\mathcal B^k$ and, for the dual complex, by $\mathcal Z_k^*$ and $\mathcal B_k^*$.}{用 $\mathcal Z^k$ 和 $\mathcal B^k$ 表示闭上链和上边界链, 对偶复形的相应空间则记为 $\mathcal Z_k^*$ 和 $\mathcal B_k^*$.}

% block: hilbert-summary-laplacian page:60 type:ordinary-paragraph
\BT{The \emph{Hodge Laplacian} is the unbounded operator $L^k$ from $W^k\to W^k$ given by $d^*d+dd^*$ with the natural domain. Its kernel is the space $\mathsf H^k=\mathcal Z^k\cap\mathcal Z_k^*$ of \emph{harmonic forms}. For a closed Hilbert complex, $\mathsf H^k$ is isomorphic to the cohomology space $\mathcal H^k$.}{\emph{Hodge Laplacian} 是从 $W^k\to W^k$ 的无界算子 $L^k$, 由 $d^*d+dd^*$ 给出并取自然定义域. 其核为\emph{调和形式}空间 $\mathsf H^k=\mathcal Z^k\cap\mathcal Z_k^*$. 对闭 Hilbert 复形, $\mathsf H^k$ 同构于上同调空间 $\mathcal H^k$.}

% block: hilbert-summary-decomposition page:60 type:ordinary-paragraph
\BT{Any closed Hilbert complex gives rise to the orthogonal \emph{Hodge decomposition}}{任一闭 Hilbert 复形都给出正交的 \emph{Hodge 分解}}
\[
W^k=\mathcal B^k\BookOrthSum\mathsf H^k\BookOrthSum\mathcal Z^k
\]
\BT{and satisfies the \emph{Poincaré inequality}}{并满足 \emph{Poincaré 不等式}}
\[
\|z\|_V\leq c_P\|dz\|,\qquad z\in\mathcal Z^{k\perp_V}.
\]

% block: hilbert-summary-solvability page:60 type:ordinary-paragraph
\BT{The Hodge Laplace problem}{Hodge Laplace 问题}
\[
L^ku=f-P_{\mathsf H}f,\qquad u\perp\mathsf H^k,
\]
\BT{has a unique solution in $D(L^k)$ for any $f\in W^k$. The solution may be characterized as well by the primal weak formulation \eqref{eq:hodge-primal-formulation} or the mixed weak formulation \eqref{eq:hodge-mixed-formulation}.}{对任意 $f\in W^k$ 在 $D(L^k)$ 中具有唯一解. 此解也可以用原始弱表述 \eqref{eq:hodge-primal-formulation} 或混合弱表述 \eqref{eq:hodge-mixed-formulation} 来刻画.}
\end{tcolorbox}
